% !TeX program = lualatex
% !TeX root = projekt 2.tex
% ============================================================
% Zajednički preamble za sva predavanja RWA kolegija
% Uključiti s: % ============================================================
% Zajednički preamble za sva predavanja RWA kolegija
% Uključiti s: % ============================================================
% Zajednički preamble za sva predavanja RWA kolegija
% Uključiti s: \input{../common/preamble.tex}
% ============================================================

\documentclass[aspectratio=149,11pt]{beamer}
\usetheme{CambridgeUS}

\usepackage{fontspec}              % radi samo sa XeLaTeX/LuaLaTeX
\setsansfont{Latin Modern Sans}    % sans-serif font
\setmonofont{Latin Modern Mono}    % monospace font za kod
\usepackage[croatian]{babel}
\usepackage{csquotes}              % \enquote{} → hrvatski navodnici „..."
\usepackage{lmodern}
\usepackage{booktabs}
\usepackage{listings}
\usepackage{xcolor}
\usepackage{hyperref}

% ---------- Autor i logo ----------
\author{mag.ing. Josip Torić (SIT UNIZD)}
\titlegraphic{\includegraphics[height=2.5cm]{../common/SIT_logo_tekst1.pdf}}

% ---------- Listings ----------
\definecolor{codebg}{RGB}{248,248,248}
\definecolor{codeframeclr}{RGB}{220,220,220}
\lstdefinelanguage{bash}{
  morekeywords={sudo,apt,brew,pip,pipx,uv,winget,python,git,flask,source,deactivate,code,gunicorn,docker,compose,heroku,render,railway,node,npm,pnpm,yarn,uvicorn,alembic,pytest,ruff,black,precommit},
  sensitive=false,
  morecomment=[l]{\#}
}
\lstdefinelanguage{python}{
  morekeywords={from,import,def,class,return,if,elif,else,for,in,while,try,except,raise,with,as,not,and,or,is,lambda,yield,True,False,None,await,async},
  sensitive=true,
  morecomment=[l]{\#},
  morestring=[b]',
  morestring=[b]"
}
\lstdefinelanguage{sql}{
  morekeywords={SELECT,FROM,WHERE,INSERT,INTO,VALUES,UPDATE,SET,DELETE,CREATE,TABLE,PRIMARY,KEY,FOREIGN,REFERENCES,UNIQUE,NOT,NULL,INDEX,ON,ALTER,DROP,CONSTRAINT},
  sensitive=false,
  morecomment=[l]{--}
}
\lstdefinelanguage{json}{
  morecomment=[l]{//},
  morestring=[b]",
  showstringspaces=false
}
\lstdefinelanguage{typescript}{
  morekeywords={import,export,from,const,let,var,function,return,if,else,for,of,in,while,switch,case,break,async,await,try,catch,finally,throw,new,class,this,extends,implements,interface,type,enum,readonly,public,private,protected,static,get,set,true,false,null,undefined,void,string,number,boolean,any,unknown,never,typeof,keyof,as},
  sensitive=true,
  morecomment=[l]{//},
  morecomment=[s]{/*}{*/},
  morestring=[b]',
  morestring=[b]",
  morestring=[b]`
}
\lstalias{javascript}{typescript}
\lstset{
  backgroundcolor=\color{codebg},
  basicstyle=\ttfamily\footnotesize,
  frame=single,
  rulecolor=\color{codeframeclr},
  frameround=tttt,
  breaklines=true,
  tabsize=2,
  showstringspaces=false,
  upquote=true,
  columns=fullflexible
}

% ---------- Helper macros ----------
\newcommand{\codefile}[1]{\texttt{#1}}
\newcommand{\kbd}[1]{\texttt{#1}}

% ---------- Code frame environment ----------
% Automatski dodaje fragile i shrink=5, koristi se umjesto \begin{frame}
% Primjer: \begin{codeframe}{Naslov} ... \end{codeframe}
% Primjer s opcijama: \begin{codeframe}[shrink=15]{Naslov} ... \end{codeframe}
\newenvironment{codeframe}[2][]
  {\begin{frame}[fragile,shrink=5,#1]{#2}}
  {\end{frame}}

% ============================================================

\documentclass[aspectratio=149,11pt]{beamer}
\usetheme{CambridgeUS}

\usepackage{fontspec}              % radi samo sa XeLaTeX/LuaLaTeX
\setsansfont{Latin Modern Sans}    % sans-serif font
\setmonofont{Latin Modern Mono}    % monospace font za kod
\usepackage[croatian]{babel}
\usepackage{csquotes}              % \enquote{} → hrvatski navodnici „..."
\usepackage{lmodern}
\usepackage{booktabs}
\usepackage{listings}
\usepackage{xcolor}
\usepackage{hyperref}

% ---------- Autor i logo ----------
\author{mag.ing. Josip Torić (SIT UNIZD)}
\titlegraphic{\includegraphics[height=2.5cm]{../common/SIT_logo_tekst1.pdf}}

% ---------- Listings ----------
\definecolor{codebg}{RGB}{248,248,248}
\definecolor{codeframeclr}{RGB}{220,220,220}
\lstdefinelanguage{bash}{
  morekeywords={sudo,apt,brew,pip,pipx,uv,winget,python,git,flask,source,deactivate,code,gunicorn,docker,compose,heroku,render,railway,node,npm,pnpm,yarn,uvicorn,alembic,pytest,ruff,black,precommit},
  sensitive=false,
  morecomment=[l]{\#}
}
\lstdefinelanguage{python}{
  morekeywords={from,import,def,class,return,if,elif,else,for,in,while,try,except,raise,with,as,not,and,or,is,lambda,yield,True,False,None,await,async},
  sensitive=true,
  morecomment=[l]{\#},
  morestring=[b]',
  morestring=[b]"
}
\lstdefinelanguage{sql}{
  morekeywords={SELECT,FROM,WHERE,INSERT,INTO,VALUES,UPDATE,SET,DELETE,CREATE,TABLE,PRIMARY,KEY,FOREIGN,REFERENCES,UNIQUE,NOT,NULL,INDEX,ON,ALTER,DROP,CONSTRAINT},
  sensitive=false,
  morecomment=[l]{--}
}
\lstdefinelanguage{json}{
  morecomment=[l]{//},
  morestring=[b]",
  showstringspaces=false
}
\lstdefinelanguage{typescript}{
  morekeywords={import,export,from,const,let,var,function,return,if,else,for,of,in,while,switch,case,break,async,await,try,catch,finally,throw,new,class,this,extends,implements,interface,type,enum,readonly,public,private,protected,static,get,set,true,false,null,undefined,void,string,number,boolean,any,unknown,never,typeof,keyof,as},
  sensitive=true,
  morecomment=[l]{//},
  morecomment=[s]{/*}{*/},
  morestring=[b]',
  morestring=[b]",
  morestring=[b]`
}
\lstalias{javascript}{typescript}
\lstset{
  backgroundcolor=\color{codebg},
  basicstyle=\ttfamily\footnotesize,
  frame=single,
  rulecolor=\color{codeframeclr},
  frameround=tttt,
  breaklines=true,
  tabsize=2,
  showstringspaces=false,
  upquote=true,
  columns=fullflexible
}

% ---------- Helper macros ----------
\newcommand{\codefile}[1]{\texttt{#1}}
\newcommand{\kbd}[1]{\texttt{#1}}

% ---------- Code frame environment ----------
% Automatski dodaje fragile i shrink=5, koristi se umjesto \begin{frame}
% Primjer: \begin{codeframe}{Naslov} ... \end{codeframe}
% Primjer s opcijama: \begin{codeframe}[shrink=15]{Naslov} ... \end{codeframe}
\newenvironment{codeframe}[2][]
  {\begin{frame}[fragile,shrink=5,#1]{#2}}
  {\end{frame}}

% ============================================================

\documentclass[aspectratio=149,11pt]{beamer}
\usetheme{CambridgeUS}

\usepackage{fontspec}              % radi samo sa XeLaTeX/LuaLaTeX
\setsansfont{Latin Modern Sans}    % sans-serif font
\setmonofont{Latin Modern Mono}    % monospace font za kod
\usepackage[croatian]{babel}
\usepackage{csquotes}              % \enquote{} → hrvatski navodnici „..."
\usepackage{lmodern}
\usepackage{booktabs}
\usepackage{listings}
\usepackage{xcolor}
\usepackage{hyperref}

% ---------- Autor i logo ----------
\author{mag.ing. Josip Torić (SIT UNIZD)}
\titlegraphic{\includegraphics[height=2.5cm]{../common/SIT_logo_tekst1.pdf}}

% ---------- Listings ----------
\definecolor{codebg}{RGB}{248,248,248}
\definecolor{codeframeclr}{RGB}{220,220,220}
\lstdefinelanguage{bash}{
  morekeywords={sudo,apt,brew,pip,pipx,uv,winget,python,git,flask,source,deactivate,code,gunicorn,docker,compose,heroku,render,railway,node,npm,pnpm,yarn,uvicorn,alembic,pytest,ruff,black,precommit},
  sensitive=false,
  morecomment=[l]{\#}
}
\lstdefinelanguage{python}{
  morekeywords={from,import,def,class,return,if,elif,else,for,in,while,try,except,raise,with,as,not,and,or,is,lambda,yield,True,False,None,await,async},
  sensitive=true,
  morecomment=[l]{\#},
  morestring=[b]',
  morestring=[b]"
}
\lstdefinelanguage{sql}{
  morekeywords={SELECT,FROM,WHERE,INSERT,INTO,VALUES,UPDATE,SET,DELETE,CREATE,TABLE,PRIMARY,KEY,FOREIGN,REFERENCES,UNIQUE,NOT,NULL,INDEX,ON,ALTER,DROP,CONSTRAINT},
  sensitive=false,
  morecomment=[l]{--}
}
\lstdefinelanguage{json}{
  morecomment=[l]{//},
  morestring=[b]",
  showstringspaces=false
}
\lstdefinelanguage{typescript}{
  morekeywords={import,export,from,const,let,var,function,return,if,else,for,of,in,while,switch,case,break,async,await,try,catch,finally,throw,new,class,this,extends,implements,interface,type,enum,readonly,public,private,protected,static,get,set,true,false,null,undefined,void,string,number,boolean,any,unknown,never,typeof,keyof,as},
  sensitive=true,
  morecomment=[l]{//},
  morecomment=[s]{/*}{*/},
  morestring=[b]',
  morestring=[b]",
  morestring=[b]`
}
\lstalias{javascript}{typescript}
\lstset{
  backgroundcolor=\color{codebg},
  basicstyle=\ttfamily\footnotesize,
  frame=single,
  rulecolor=\color{codeframeclr},
  frameround=tttt,
  breaklines=true,
  tabsize=2,
  showstringspaces=false,
  upquote=true,
  columns=fullflexible
}

% ---------- Helper macros ----------
\newcommand{\codefile}[1]{\texttt{#1}}
\newcommand{\kbd}[1]{\texttt{#1}}

% ---------- Code frame environment ----------
% Automatski dodaje fragile i shrink=5, koristi se umjesto \begin{frame}
% Primjer: \begin{codeframe}{Naslov} ... \end{codeframe}
% Primjer s opcijama: \begin{codeframe}[shrink=15]{Naslov} ... \end{codeframe}
\newenvironment{codeframe}[2][]
  {\begin{frame}[fragile,shrink=5,#1]{#2}}
  {\end{frame}}


\title[Programiranje za web]{Programiranje za web \\ Projekt 2 - Frontend}
\date{23.11.2026.}

\begin{document}

\frame{\titlepage}

% ==============================
% DIO 1: KONTEKST I TEHNOLOGIJE
% ==============================

% ------------------------------
\begin{frame}{Plan prezentacije}
\begin{enumerate}
  \item Što se traži u 2. dijelu projekta
  \item Tehnologije i instalacija
  \item Detaljni zahtjevi za frontend
  \item Bodovna raspodjela 2. dijela
  \item Što vas čeka na 2. obrani
\end{enumerate}
\end{frame}

% ------------------------------
\begin{frame}{Gdje smo: backend je obranjen}
Iza vas je 1. obrana - imate funkcionalan API s autentikacijom, autorizacijom i poslovnim pravilima. Sada gradite \textbf{SPA frontend koji komunicira s vašim API-jem}.

\vspace{0.6em}
\begin{tabular}{@{}lp{5.2cm}lr@{}}
  \toprule
  \textbf{Dio} & \textbf{Što se predaje} & \textbf{Obrana} & \textbf{Bodovi} \\
  \midrule
  1. dio & Backend: API, baza, testovi & 16.11. & 45 \\
  2. dio & Frontend SPA + deploy & 18.01. & 45 \\
  \multicolumn{2}{@{}l}{Dolasci na predavanja} & & 10 \\
  \midrule
  \multicolumn{3}{@{}l}{\textbf{Ukupno}} & \textbf{100} \\
  \bottomrule
\end{tabular}

\vspace{0.6em}
Frontend se predaje kao \codefile{web/} folder u istom monorepou, a cijela aplikacija mora biti \textbf{deployana na javni URL}.
\end{frame}

% ------------------------------
\begin{frame}{Tehnologije}
\begin{tabular}{@{}lp{7cm}@{}}
  \toprule
  \textbf{Tehnologija} & \textbf{Uloga} \\
  \midrule
  Vue 3 & UI framework, Composition API \\
  TypeScript & jezik frontenda \\
  Vite & dev server i build tool \\
  Vue Router & routing i guardovi \\
  Pinia & state management \\
  Axios & HTTP komunikacija s API-jem \\
  Vitest + @vue/test-utils & unit i component testovi \\
  Playwright & E2E testovi \\
  Railway ili Render & deploy \\
  \bottomrule
\end{tabular}

\vspace{0.4em}
Kao i za backend: drugi stack nije dopušten. React, Angular ili čisti JavaScript nisu opcija.
\end{frame}

% ------------------------------
\begin{frame}[fragile,shrink=5]{Instalacija: Node.js i ekstenzije}
\begin{lstlisting}[language=bash]
# Node.js - preuzmite LTS sa https://nodejs.org
node --version   # 22+
npm --version

# Scaffold projekta (radimo zajedno na predavanju)
npm create vue@latest
# odaberite: TypeScript = Yes, Router = Yes, Pinia = Yes
\end{lstlisting}

\vspace{0.6em}
Preporučene VS Code ekstenzije za frontend:
\begin{itemize}
  \item \textbf{Vue - Official} - syntax highlighting, TypeScript podrška u .vue datotekama
  \item \textbf{ESLint} + \textbf{Prettier} - lint i formatiranje
\end{itemize}
\end{frame}

% ------------------------------
\begin{frame}{Struktura frontenda}
\begin{tabular}{@{}lp{7.5cm}@{}}
  \toprule
  \textbf{Folder} & \textbf{Sadržaj} \\
  \midrule
  \codefile{web/src/views/} & stranice - jedna po ruti \\
  \codefile{web/src/components/} & reusable UI komponente \\
  \codefile{web/src/stores/} & Pinia storeovi \\
  \codefile{web/src/router/} & rute i guardovi \\
  \codefile{web/src/api/} & Axios client i API pozivi \\
  \codefile{web/src/types/} & TypeScript tipovi, 1:1 s Pydantic schemama \\
  \codefile{web/src/utils/} & pomoćne funkcije \\
  \bottomrule
\end{tabular}

\vspace{0.4em}
Ista disciplina kao na backendu: svaki folder ima jednu odgovornost. API pozivi ne žive po komponentama, nego u \codefile{api/} sloju.
\end{frame}

% ==============================
% DIO 2: DETALJNI ZAHTJEVI
% ==============================

% ------------------------------
\begin{frame}{Zahtjev 1: SPA s routingom}
\begin{itemize}
  \item \textbf{Najmanje 5 stranica}: login, lista glavnog entiteta, forma za unos ili uređivanje, detalji, admin pregled
  \item Vue Router s \codefile{createWebHistory} - navigacija bez refresha stranice
  \item Zajednički \textbf{layout}: navigacija s linkovima ovisnim o roli, logout gumb, prikaz prijavljenog korisnika
  \item Dinamičke rute s parametrima: \codefile{/lifters/:id/edit} ili ekvivalent u vašoj domeni
\end{itemize}

\vspace{0.4em}
Na obrani otvaram DevTools i klikćem po aplikaciji: ako se stranica refresha na svaku navigaciju, to nije SPA.
\end{frame}

% ------------------------------
\begin{frame}{Zahtjev 2: auth flow}
\begin{itemize}
  \item \textbf{Login forma} koja poziva vaš \codefile{/auth/login} i sprema tokene
  \item Access token se šalje na svaki request kroz \textbf{Axios request interceptor}
  \item \textbf{Automatski refresh}: na 401 interceptor tiho zove \codefile{/auth/refresh} i ponavlja originalni request
  \item Kad refresh ne uspije: logout i redirect na login
  \item Auth stanje preživljava refresh stranice - \codefile{/auth/me} pri startu aplikacije
\end{itemize}

\vspace{0.4em}
Ovo testiram tako da ostavim aplikaciju otvorenu dok access token ne istekne - korisnik ne smije primijetiti ništa.
\end{frame}

% ------------------------------
\begin{frame}{Zahtjev 3: route guardovi i permission-aware UI}
\begin{itemize}
  \item \codefile{router.beforeEach} guard: neprijavljeni korisnik ide na login
  \item Admin rute nedostupne običnom korisniku - direktan unos URL-a preusmjerava
  \item UI poštuje rolu: korisnik ne vidi admin akcije, admin vidi sve
  \item UI poštuje \textbf{poslovna pravila}: gumbi disabled kad pravilo zabranjuje akciju, npr. nakon isteka roka
\end{itemize}

\vspace{0.6em}
Guard je UX, ne sigurnost - backend i dalje mora zabraniti. Očekujte pitanje: \enquote{Što se dogodi ako korisnik ručno upiše admin URL? A ako pošalje request curlom?}
\end{frame}

% ------------------------------
\begin{frame}{Zahtjev 4: UX stanja}
Svaki view koji dohvaća podatke mora hendlati tri stanja:

\vspace{0.4em}
\begin{tabular}{@{}llp{5.5cm}@{}}
  \toprule
  \textbf{Stanje} & \textbf{Uvjet} & \textbf{Što korisnik vidi} \\
  \midrule
  Loading & čeka se API & spinner ili skeleton \\
  Error & request pao & poruka s retry gumbom \\
  Empty & nema podataka & jasna poruka i CTA za prvi unos \\
  \bottomrule
\end{tabular}

\vspace{0.6em}
Prazna bijela stranica dok se čeka API, ili tablica bez retka i bez objašnjenja - to su stanja za koja skidam bodove. Retry nudite samo za mrežne i 5xx greške, ne za 4xx.
\end{frame}

% ------------------------------
\begin{frame}{Zahtjev 5: forme i validacija}
\begin{itemize}
  \item \textbf{Klijentska validacija} prije slanja: obavezna polja, formati, rasponi
  \item \textbf{Serverske greške} prikazane na pravom mjestu:
  \begin{itemize}
    \item 422 - greška ispod odgovarajućeg polja, parsirana iz FastAPI \codefile{detail} formata
    \item 409 i 400 - globalna poruka iznad forme, npr. \enquote{Rok za prijave je istekao}
  \end{itemize}
  \item Submit gumb disabled dok request traje - nema duplih klikova
  \item Nakon uspjeha: redirect ili jasna potvrda, ne tiho ništa
\end{itemize}

\vspace{0.4em}
Pri pregledu namjerno šaljem nevažeće podatke i duplikat - forma mora reagirati smisleno na oba slučaja.
\end{frame}

% ------------------------------
\begin{frame}{Zahtjev 6: state management}
\begin{itemize}
  \item \textbf{Pinia auth store}: user, token, \codefile{isAuthenticated}, \codefile{isAdmin}, login i logout akcije
  \item \textbf{Barem jedan store za poslovni entitet} s loading i error stateom
  \item Store za podatke koje dijeli više viewova; lokalni \codefile{ref()} za ono što živi u jednoj komponenti
  \item Pesimistički update: čekaj odgovor API-ja, pa ažuriraj state
\end{itemize}

\vspace{0.4em}
Pitanje koje sigurno dolazi: \enquote{Zašto je ovo u storeu, a ono lokalno?} Odgovor \enquote{tako je bilo u tutorialu} ne prolazi.
\end{frame}

% ------------------------------
\begin{frame}{Zahtjev 7: TypeScript}
\begin{itemize}
  \item Cijeli frontend u TypeScriptu: \codefile{lang="ts"} u svim komponentama
  \item \textbf{Tipovi u \codefile{src/types/}} preslikavaju vaše Pydantic Response scheme 1:1
  \item Tipizirani API pozivi: \codefile{api.get<Lifter[]>(...)}, tipizirani propsi i emitovi
  \item \codefile{vue-tsc} prolazi \textbf{bez grešaka} - to je uvjet, ne preporuka
  \item \codefile{any} koristite samo gdje imate obrazloženje koje možete izgovoriti naglas
\end{itemize}

\vspace{0.4em}
TypeScript vam je saveznik protiv tipfelera u nazivima polja - greška se vidi u editoru, ne na obrani.
\end{frame}

% ------------------------------
\begin{frame}{Zahtjev 8: testovi}
\begin{tabular}{@{}lp{6.8cm}@{}}
  \toprule
  \textbf{Tip} & \textbf{Minimum} \\
  \midrule
  Vitest unit & barem jedna utility funkcija i jedan store s mockanim API-jem \\
  Component test & barem jedna komponenta: renderiranje propsa i emitani event \\
  Playwright E2E & barem login flow: uspješna prijava i pogrešna lozinka \\
  \bottomrule
\end{tabular}

\vspace{0.6em}
Testovi se pokreću s \codefile{npm run test} i moraju biti zeleni pri predaji. E2E test poslovne zabrane, npr. disabled gumb nakon roka, nosi dodatnu težinu na obrani.
\end{frame}

% ------------------------------
\begin{frame}{Zahtjev 9: deploy}
\begin{itemize}
  \item Aplikacija dostupna na \textbf{javnom URL-u}: Railway, Render ili slična platforma
  \item Backend i frontend deployani zasebno; baza je managed PostgreSQL
  \item \textbf{Env varijable}: \codefile{VITE\_API\_URL} za frontend; \codefile{DATABASE\_URL}, \codefile{JWT\_SECRET} i \codefile{CORS\_ORIGINS} za backend
  \item \textbf{CORS} ograničen na domenu vašeg frontenda, ne \codefile{*}
  \item \textbf{SPA fallback}: direktan URL bilo koje rute servira \codefile{index.html}
  \item Migracije se pokreću automatski pri deployu
\end{itemize}

\vspace{0.4em}
Deploy radimo zajedno na Predavanju 12 - ne ostavljajte ga za zadnji dan.
\end{frame}

% ------------------------------
\begin{frame}{Bodovna raspodjela 2. dijela (45 bodova)}
\begin{tabular}{@{}lr@{}}
  \toprule
  \textbf{Komponenta} & \textbf{Bodovi} \\
  \midrule
  Routing, layout i guardovi & 6 \\
  Auth flow: login, refresh, 401 handling & 7 \\
  CRUD viewovi s UX stanjima & 8 \\
  Forme i prikaz serverskih grešaka & 5 \\
  State management (Pinia) & 5 \\
  TypeScript kvaliteta & 4 \\
  Testovi: Vitest + Playwright & 5 \\
  Deploy & 5 \\
  \midrule
  \textbf{Ukupno} & \textbf{45} \\
  \bottomrule
\end{tabular}

\vspace{0.4em}
Isto pravilo kao u 1. dijelu: dio koji na usmenom ne znate objasniti ne nosi bodove.
\end{frame}

% ------------------------------
\begin{frame}{Checklist za predaju 2. dijela}
\begin{enumerate}
  \item \codefile{npm install} i \codefile{npm run dev} rade na svježem cloneu
  \item \codefile{vue-tsc} prolazi bez grešaka
  \item \codefile{npm run test} je zelen
  \item \codefile{npm run build} prolazi bez grešaka
  \item Javni URL radi: login, CRUD, pravila - sve na deployanoj verziji
  \item README dopunjen: pokretanje frontenda i link na deploy
\end{enumerate}

\vspace{0.4em}
Na obrani radimo na \textbf{deployanoj} aplikaciji, ne na localhost verziji. Ako deploy ne radi, obrana staje na tome.
\end{frame}

% ==============================
% DIO 3: OBRANA
% ==============================

% ------------------------------
\begin{frame}{AI alati: ista pravila kao u 1. dijelu}
Sve iz prezentacije \enquote{Projekt 1} i dalje vrijedi: AI smijete koristiti, ali na obrani morate razumjeti i objasniti svaki dio svog koda.

\vspace{0.6em}
Za frontend posebno: AI rado generira komponente koje rade, ali s \codefile{any} tipovima, logikom u templateu i copy-paste stilom. Prođite kroz generirano i dovedite ga na razinu zahtjeva - to je dio posla, ne opcija.
\end{frame}

% ------------------------------
\begin{frame}{2. obrana (18.01.): kako izgleda}
Demonstracija deployane aplikacije, pregled koda, pitanja.

\vspace{0.4em}
\textbf{Tipična pitanja:}
\begin{itemize}
  \item Kako frontend hendla 401 odgovor? Gdje je interceptor?
  \item Što se dogodi kad korisnik ručno upiše URL admin stranice?
  \item Gdje držiš stanje i zašto baš tu?
  \item Što se dogodi kad token istekne dok korisnik radi?
  \item Kako si konfigurirao CORS i koje env varijable si postavio?
  \item Što testira ovaj Vitest test? Što pokriva E2E, a unit ne može?
\end{itemize}

\vspace{0.4em}
Na kraju: \textbf{mini promjena uživo} - dodaj polje u formu, prikaži novi podatak, prilagodi tip.
\end{frame}

% ------------------------------
\begin{frame}{Što je dovoljan odgovor na usmenom}
\begin{columns}[T]
\column{0.48\textwidth}
\textbf{Prolazi}
\begin{itemize}
  \item \enquote{Interceptor uhvati 401, pozove refresh, ponovi originalni request s \_retry flagom}
  \item \enquote{Guard čita meta.requiresAdmin i rolu iz storea pa preusmjerava}
  \item \enquote{Store jer listu koriste dva viewa - forma je lokalni state}
\end{itemize}

\column{0.48\textwidth}
\textbf{Ne prolazi}
\begin{itemize}
  \item \enquote{To mi je AI napisao, ne znam točno}
  \item \enquote{Mislim da router to nekako sredi}
  \item \enquote{Radi na localhostu, ne znam zašto na deployu ne radi}
\end{itemize}
\end{columns}
\end{frame}

% ------------------------------
\begin{frame}{Rokovi i sljedeći koraci}
\begin{tabular}{@{}ll@{}}
  \toprule
  \textbf{Što} & \textbf{Kada} \\
  \midrule
  Scaffold frontenda u \codefile{web/} & ovaj tjedan \\
  Predavanja 7-11: frontend gradivo & 23.11. - 21.12. \\
  Predavanje 12: deploy zajedno & 11.01. \\
  2. obrana: full-stack na javnom URL-u & 18.01. \\
  Rezervni termin predaje & veljački rokovi \\
  \bottomrule
\end{tabular}

\vspace{0.6em}
Ritam je isti kao u 1. dijelu: svako predavanje reproducirajte na svom projektu isti tjedan. Između zadnjeg predavanja i obrane je samo tjedan dana - tko sve ostavi za praznike, ne stigne.

\vspace{0.6em}
\begin{center}
\Large Pitanja?
\end{center}
\end{frame}

\end{document}
